% ==============================================================
%  Chapter 6: AI for Teachers and Educators
% ==============================================================
\chapter{AI for Teachers and Educators}
\label{ch:teachers}

Teaching is fundamentally about understanding what the learner needs and adapting
delivery accordingly.
AI can dramatically reduce the time spent on material production, personalisation, and
administrative tasks --- freeing teachers for the high-value work that requires human
presence: mentoring, facilitating discussion, and building relationships.

\section{Curriculum Design and Content Creation}
\label{sec:curriculum}

\textbf{What AI does well:}
\begin{itemize}
  \item Generating lesson outlines from learning objectives
  \item Producing multiple explanations of the same concept at different complexity levels
  \item Creating worked examples, analogies, and counterexamples
  \item Drafting slides, handouts, and reading materials
  \item Translating technical content into age-appropriate language
\end{itemize}

\textbf{Effective prompts:}
\begin{itemize}
  \item \textit{``Create a 90-minute lesson plan on recursion for second-year computer science students. Include: learning objectives, a real-world hook, a worked example, a group activity, and 3 check-for-understanding questions.''}
  \item \textit{``Explain the concept of compound interest three ways: for a 12-year-old, for a high school student, and for an economics undergraduate.''}
  \item \textit{``Generate 5 common student misconceptions about [topic] and how to address each.''}
\end{itemize}

\begin{tipbox}
Ask AI to generate the misconceptions list for a topic before you teach it.
\textit{``What do students typically misunderstand about [concept]?''}
This prepares you to pre-empt confusion rather than react to it.
\end{tipbox}

\section{Personalised Learning and Differentiation}
\label{sec:personalisation}

AI allows teachers to create differentiated materials at scale --- something that was
previously prohibitively time-consuming.

\begin{itemize}
  \item \textit{``Rewrite this explanation for a student with dyslexia: shorter sentences, concrete examples, no abstract metaphors.''}
  \item \textit{``Create an extension activity for fast-finishers that deepens understanding of [concept] rather than just adding more practice.''}
  \item \textit{``Simplify this problem set for students who are still developing [prerequisite skill].''}
\end{itemize}

\section{Assessment and Feedback}
\label{sec:assessment}

\textbf{Assessment design:}
\begin{itemize}
  \item \textit{``Write 10 multiple-choice questions on [topic] at Bloom's levels: 3 recall, 4 application, 3 analysis. Include a worked solution for each.''}
  \item \textit{``Design a project rubric for [assignment] with four performance levels across these criteria: [list].''}
\end{itemize}

\textbf{Feedback drafting:}
AI can draft structured feedback on student work when given a rubric and the student
submission.
Paste the rubric, the student's work, and ask:
\textit{``Draft constructive feedback highlighting two strengths and two specific
improvements, referenced to the criteria.''}
You then personalise the tone and verify accuracy.

\begin{warnbox}
AI-generated feedback is a starting point, not a final product.
Students need feedback that reflects genuine engagement with their work.
Review and personalise every AI-drafted response before sending.
\end{warnbox}

\section{Ethical Considerations in Education}
\label{sec:edu-ethics}

The challenge AI presents in education is not primarily about plagiarism detection ---
it is about ensuring students actually learn.

\textbf{Key principles:}
\begin{itemize}
  \item Design assessments that require personal experience, live demonstration, or
        iterative conversation --- these are hard for AI to fake well.
  \item Teach students \emph{how} to use AI effectively rather than banning it.
        They will use it regardless; they might as well learn to use it well.
  \item Use AI detection tools as one signal, not a verdict.
        False positives are real and damaging.
  \item Model appropriate AI use yourself: show students when you used AI and how
        you verified and edited its output.
\end{itemize}

\begin{infobox}[Reframing the Question]
The question is not \textit{``did the student use AI?''} but \textit{``can the student
demonstrate understanding and apply the concept independently?''}
Design your assessments to answer the second question.
\end{infobox}
